%% kapa: exam=1 round=29 year=2018 period=1 form=A subject=civil_law subject_ko=민법 seq=14 number=14
%% source: Q-Net 감정평가사 기출문제 (cst00309) artlSeq=5202808 file=extracted/1008/5202808/2018년 제29회 감정평가사 1차시험 기출문제/2018년 제29회 감정평가사 1차시험 기출문제/1교시 A형.pdf
%% conversion: kordoc 4.24.51 -> tools/gen1.py (auto); review=proofread-sonnet; flags=none
\begin{question}{14}
법률행위의 무효에 관한 설명으로 옳지 않은 것은? (다툼이 있으면 판례에 따름)
\begin{choices}
  \item 강박의 정도가 극심하여 의사결정을 스스로 할 수 있는 여지가 완전히 박탈된 상태에서 의사표시가 이루어진 경우 그 의사표시는 무효이다.
  \item 반사회적 법률행위를 원인으로 부동산에 관한 소유권이전등기를 마친 등기명의자가 소유권에 기한 물권적 청구권을 행사하는 경우, 상대방은 법률행위의 무효를 항변으로서 주장할 수 없다.
  \item 무효인 법률행위를 추인에 의하여 새로운 법률행위로 보기 위해서는 당사자가 이전의 법률행위가 무효임을 알고 그 행위에 대하여 추인하여야 한다.
  \item 무효인 법률행위가 다른 법률행위의 요건을 구비하고 당사자가 그 무효를 알았더라면 다른 법률행위를 하는 것을 의욕하였으리라고 인정될 때에는 다른 법률행위로서 효력을 가진다.
  \item 후속행위를 한 것이 묵시적 추인으로 인정되기 위해서는 이전의 법률행위가 무효임을 알거나, 무효임을 의심하면서도 그 행위의 효과를 자기에게 귀속시키도록 하는 의사로 후속행위를 하였음이 인정되어야 한다.
\end{choices}
\end{question}
\answer{2}
